% Conteggio requisiti:
% Funzionali Obbligatori: R1..R47-F-O, R48..R51-F-O, R53..R77-F-O = 76 requisiti
% Funzionali Desiderabili: R52-F-D = 1 requisito
% Funzionali Opzionali: R78..R86-F-Op = 9 requisiti
% Qualità Obbligatori: R1-Q-O, R2-Q-O, R3-Q-O, R4-Q-O, R5-Q-O, R6-Q-O = 6 requisiti
% Qualità Desiderabili: R9-Q-D = 1 requisito
% Qualità Opzionali: R7-Q-Op, R8-Q-Op = 2 requisiti
% Vincolo Obbligatori: R1-V-O, R2-V-O, R3-V-O, R5-V-O, R6-V-O, R7-V-O = 6 requisiti
% Vincolo Opzionali: R4-V-Op = 1 requisito
% Prestazionali Obbligatori: R1-P-O, R2-P-O = 2 requisiti
% Totale Obbligatori: 76 + 6 + 6 + 2 = 90
% Totale Desiderabili: 1 + 1 + 0 + 0 = 2
% Totale Opzionali: 9 + 2 + 1 + 0 = 12

\section{Riepilogo}
Nella seguente sezione viene riportato un riepilogo quantitativo dello stato dei requisiti individuati, suddivisi in base al loro grado di importanza.

\subsection{Requisiti obbligatori individuati}
\begin{table}[ht]
	\centering
	\renewcommand{\arraystretch}{1.4}
	\begin{tabularx}{0.8\textwidth}{|l|X|}
		\hline
		\rowcolor{primarycolor!10}
		\textbf{Tipologia Requisito} & \textbf{Numero Requisiti Obbligatori} \\
		\hline
		Funzionali & 76 \\
		\hline
		Di Qualità & 6 \\
		\hline
		Di Vincolo & 6 \\
		\hline
		Prestazionali & 2 \\
		\hline
		\textbf{Totale Obbligatori} & \textbf{90} \\
		\hline
	\end{tabularx}
	\caption{Riepilogo requisiti obbligatori}
\end{table}

\subsection{Requisiti desiderabili individuati}
\begin{table}[ht]
	\centering
	\renewcommand{\arraystretch}{1.4}
	\begin{tabularx}{0.8\textwidth}{|l|X|}
		\hline
		\rowcolor{primarycolor!10}
		\textbf{Tipologia Requisito} & \textbf{Numero Requisiti Desiderabili} \\
		\hline
		Funzionali & 1 \\
		\hline
		Di Qualità & 1 \\
		\hline
		Di Vincolo & 0 \\
		\hline
		Prestazionali & 0 \\
		\hline
		\textbf{Totale Desiderabili} & \textbf{2} \\
		\hline
	\end{tabularx}
	\caption{Riepilogo requisiti desiderabili}
\end{table}

\subsection{Requisiti opzionali individuati}
\begin{table}[ht]
	\centering
	\renewcommand{\arraystretch}{1.4}
	\begin{tabularx}{0.8\textwidth}{|l|X|}
		\hline
		\rowcolor{primarycolor!10}
		\textbf{Tipologia Requisito} & \textbf{Numero Requisiti Opzionali} \\
		\hline
		Funzionali & 9 \\
		\hline
		Di Qualità & 2 \\
		\hline
		Di Vincolo & 1 \\
		\hline
		Prestazionali & 0 \\
		\hline
		\textbf{Totale Opzionali} & \textbf{12} \\
		\hline
	\end{tabularx}
	\caption{Riepilogo requisiti opzionali}
\end{table}

\subsection{Riepilogo generale}
\begin{table}[ht]
	\centering
	\renewcommand{\arraystretch}{1.4}
	\begin{tabularx}{0.8\textwidth}{|l|X|}
		\hline
		\rowcolor{primarycolor!10}
		\textbf{Grado di Importanza} & \textbf{Totale Requisiti} \\
		\hline
		Obbligatori & 90 \\
		\hline
		Desiderabili & 2 \\
		\hline
		Opzionali & 12 \\
		\hline
		\textbf{Totale complessivo} & \textbf{104} \\
		\hline
	\end{tabularx}
	\caption{Riepilogo generale dei requisiti}
\end{table}